\section{Results: complete static information remains a strong control}
\subsection{The complete-service comparison}
Table~\ref{tab:formal} reports every formal group. Average-energy yields
8 PASS and 4 FAIL. Rich-static, complete-input, executable-response and
E4 table each yield 12 PASS. Live E5 yields 7 PASS and 5 INCOMPLETE;
it has no admitted final service FAIL. Its completion rate is 7/12,
not a conditional 7/7. Missing final executions and submitted service
failures remain different outcomes with the same full denominator.

% Generated from hash-verified saved data; no scientific execution.
\begin{table}[htbp]
\centering\small
\caption{Formal matched interface comparison: four fixed tasks, three repetitions per task, twelve groups per arm. P/F/I are PASS, submitted FAIL and INCOMPLETE. Every incomplete group and its requests remain included.}
\label{tab:formal}
\begin{tabular}{@{}lrrrrr@{}}\toprule
Condition & P & F & I & Requests & Requests/group \\\midrule
Average energy & 8 & 4 & 0 & 75 & 6.2500 \\
Rich static / E4 JSON & 12 & 0 & 0 & 65 & 5.4167 \\
Complete input & 12 & 0 & 0 & 64 & 5.3333 \\
Executable response & 12 & 0 & 0 & 66 & 5.5000 \\
Live observation (E5) & 7 & 0 & 5 & 80 & 6.6667 \\
E4 typed table & 12 & 0 & 0 & 65 & 5.4167 \\
\bottomrule\end{tabular}
\end{table}

% Generated from hash-verified saved data; no scientific execution.
\begin{table}[htbp]
\centering\small
\caption{Formal P/F/I out of three per task. Each strong arm separately denotes rich static, complete input, executable response and E4 table; none has an omitted failure.}
\label{tab:tasks}
\begin{tabular}{@{}lccc@{}}\toprule
Fixed task & Average energy & Live E5 & Each strong arm \\\midrule
Pulse-fast & 3/0/0 & 2/0/1 & 3/0/0 \\
Volatile reset & 1/2/0 & 2/0/1 & 3/0/0 \\
Latched reset & 1/2/0 & 1/0/2 & 3/0/0 \\
C-lite query derivative & 3/0/0 & 2/0/1 & 3/0/0 \\
\bottomrule\end{tabular}
\end{table}


The coarse arm's observed deficit occurs in the two reset tasks, each
with one PASS and two FAIL. Pulse-fast and the C-lite derivative each
have three PASS. Every strong arm passes all three repetitions on each
task (Table~\ref{tab:tasks}). This finite result is consistent with the
value of supplied timing and state information beyond the coarse energy
summary. It establishes no service advantage for executable-response
over rich-static information. The saturated controls also do not prove
equivalence on harder or previously unused tasks.

The two pilot repetitions are not pooled with the formal twelve. On
their single volatile-reset input, average-energy has one PASS and one
FAIL; each other displayed arm has two PASS. E5's two pilot completions
include one selection and one construction. These are the development
observations preceding the one-time formal-repeat freeze, not additional
formal successes or independent tasks~\cite{tables,protocol}.

\subsection{What the seven prospective E5 chains did}
Every formal E5 completion contains an actual prospective proposal,
external source qualification and a typed observation consumed by the
design policy. Six select existing fields and one constructs a relation.
Five use component records and two excitation records. Six change the
initial design and one does not. All seven final designs pass independent
service, and all seven deliver before any reset. Their successes therefore
establish neither pending-state survival across a reset nor the necessity
of the acquired observation. New physical measurements number
zero~\cite{tables,comparison}.

% Generated from hash-verified saved data; no scientific execution.
\begin{table}[htbp]
\centering\small
\caption{Disposition of all twelve formal E5 groups. Repeated rejected actions are not additional groups; absent executions are retained as INCOMPLETE.}
\label{tab:e5}
\begin{tabular}{@{}lrc@{}}\toprule
Disposition & Groups & Final service \\\midrule
Qualified selection, consumed & 6 & PASS \\
Qualified construction, consumed & 1 & PASS \\
Source-unit binding rejected & 4 & INCOMPLETE \\
No typed consumer match & 1 & INCOMPLETE \\
\bottomrule\end{tabular}
\end{table}


Four incomplete groups retain \path{SOURCE_UNIT_UNAVAILABLE} in their
original gateway records. The code identifies a submitted source/field/domain
unit-binding problem; it does not imply that a physical quantity does not
exist. In the pulse-fast example, the selected packet field
\path{/packets/0/application_bytes} is unit-mapped in the acquired
saved-activity source, while the declared domain also references deadline,
TX-duration and readiness fields absent from that source's unit map.
A latched-reset proposal selects a packet field from the model-law source,
whose mapped fields describe model identity and Python definition.
The original generic error does not name the first rejected term.
We diagnose the source-local mismatch from saved proposals and logs;
no repaired proposal is evaluated~\cite{gateways}.

The fifth incomplete group records \path{NO_TYPED_CONSUMER_MATCH} on
the C-lite derivative. Its equality expression outputs a Boolean while
the consumer branches match numeric durations 0.002 and 0.001. A value
can therefore be syntactically meaningful yet fail to enter the intended
design edge. Other syntax and valid-proposal-cap errors remain charged
actions inside these groups, not extra independent trials or free
rescue attempts. Table~\ref{tab:e5} counts groups once.

The rejected original E5 attempt and its separately authorized post-hoc
adapter chain remain outside these twelve formal groups. The post-hoc
record demonstrates engineering linkage of an already generated relation.
The seven later formal completions demonstrate prospective operation
under a different registered protocol. Relabeling the former as a
successful live discovery would erase both its original admission
failure and the host intervention~\cite{original,posthoc}.

\subsection{Equal information and acquisition overhead}
E4 gives identical complete-service and request totals: JSON and the
typed table each pass 12/12 using 65 provider requests, or 5.4167 per
group. The retained USD and wall receipts differ
(Table~\ref{tab:format}), but this small comparison isolates neither
all latency determinants nor the representation of later acquisitions.
It supplies no observed service or request advantage and no causal
format-speed or format-cost conclusion~\cite{tables}.

% Generated from hash-verified saved data; no scientific execution.
\begin{table}[htbp]
\centering\small
\caption{E4 equal-information initial serialization. Later acquisitions use common JSON. Cost and summed group wall time are descriptive; neither establishes a format advantage.}
\label{tab:format}
\begin{tabular}{@{}lrrrr@{}}\toprule
Initial packet & PASS & Requests & USD & Wall (s) \\\midrule
JSON & 12/12 & 65 & 0.4267200 & 676.5 \\
Typed table & 12/12 & 65 & 0.5145681 & 611.4 \\
\bottomrule\end{tabular}
\end{table}

% Generated from hash-verified saved data; no scientific execution.
\begin{table}[htbp]
\centering\small
\caption{Finite observation selection on four exposed cases. Read and audit bytes are Python payloads, not disk I/O. Host timings are single historical measurements, not latency benchmarks; selected and mechanism-spec contain the same facts.}
\label{tab:acquisition}
\begin{tabular}{@{}lrrrrr@{}}\toprule
Observation & Correct & Read B & Audit B & Wall ms & Model runs \\\midrule
uninformative & 0/4 & 4648 & 4648 & 0.289 & 0 \\
mechanism-spec & 4/4 & 103 & 103 & 0.260 & 0 \\
fixed-rich & 4/4 & 4747 & 4747 & 0.282 & 0 \\
full-raw & 4/4 & 8544 & 8544 & 0.333 & 0 \\
executable-dynamic & 4/4 & 8544 & 8544 & 1.886 & 4 \\
selected & 4/4 & 103 & 103 & 0.236 & 0 \\
\bottomrule\end{tabular}
\end{table}


The earlier finite observation selector chooses the smallest eligible
source using researcher-supplied required-field metadata, without service
outcomes. It reads the same 103 bytes and facts as the fixed mechanism-spec
baseline; both correctly judge four cases. Rich-static, complete raw and
an executable response also correctly judge all four. The historical
timings in Table~\ref{tab:acquisition} include acquisition and source
audit, plus selection when used. They are single-host payload measurements
with uncontrolled cache state, not storage-I/O or monetary measurements.
Thus lower payload volume relative to a full model is not an independent
selection-algorithm gain~\cite{structural}.

A deterministic compiler also constructs the researcher-specified
headroom observation. Its saved outputs of $-14$ and $8~\mu$J agree
with twenty exposed clamp-grid checks and two fresh-process executions.
Construction took 0.239 ms, reference qualification 13.017 ms and the
two fresh processes 37.730 ms combined in those retained receipts.
Because the equation and relevant fields were supplied by a researcher,
this is executable construction machinery, not an autonomous discovery.

\subsection{Service success versus an ordered retention witness}
Four formal participant groups select a pending-event-across-reset result:
two complete-input groups, one proposed-method group and one E4 table
group. They make four fresh choices but reuse one distinct historical
digital execution key. All select $10~\mu$F and earliest TX at 45 ms.
The same event is captured from 3.7 to 7.9 ms, a brownout occurs at
12.1 ms, and first valid reception follows at 46.9 ms. The fixed ideal
passive latch provides the retention mechanism~\cite{ordered}.

All twelve rich-static, all twelve executable-response and all seven
successful E5 groups instead deliver before any reset. Their complete
service passes are valid for their tasks but are not evidence for the
stronger ordered-retention property. Conversely, four selections of one
response are not four independent physical experiments or four newly
invented retention mechanisms. The ordered audit creates zero new
scientific executions and retains physical qualification UNKNOWN.
\FloatBarrier
